\section{A discriminant denominator for harmonic derivatives}
\label{sec:geometry-discriminant-clearing}

The denominator needed for integral theta series permits a simpler local
argument than a sharp estimate for every ramified eigenspace.  We prove
the required, coarser bound directly.  The formal group, biextension,
rigidifications, and integral sections in the statement are fixed geometric
data; their construction is a separate matter.

Let $R$ be a complete discrete valuation ring of characteristic zero, with
fraction field $K$.  Let $E$ be a CM field whose embeddings take values in
$K$, put $O=\mathcal O_E$ and $\Delta=|\operatorname{disc} E|$, and write
$\bar\tau$ for the conjugate of an embedding $\tau$.  Consider a smooth
commutative formal group $G$ of dimension $g$ over $R$, with an $O$-action
and an integral formal biextension.  Retain its rigidifications, symmetry,
and adjoint identification for complex conjugation.  Let $h(x,z)$ be a
canonical logarithmic trivialization over $K$, vanishing on both axes.
Assume that the given element $w\in O$ satisfies $w+\bar w=1$, and set
\[
 h_L(x)=h(x,[w]x),\qquad
 \sigma_i(x)=-\left.\frac{\partial h(x,z)}{\partial z_i}\right|_{z=0},
 \qquad S_{ij}=\left.\frac{\partial\sigma_i}{\partial x_j}\right|_{x=0}.
\]
All derivatives below are ordinary derivatives, without divided powers.

For each embedding $\tau$, choose finitely many integral tangent vectors
$u_{\tau i}\in\operatorname{Lie}(G)(R)$ satisfying
$[a]u_{\tau i}=\tau(a)u_{\tau i}$ for $a\in O$.
For $n\geq0$, let $P$ be a polynomial with coefficients in $R$ in the
variables $X_{\tau i b}$, $1\leq b\leq n$, homogeneous of total degree $d$.
Its required harmonicity is
\begin{equation}\label{eq:discriminant-harmonic}
 \sum_{b=1}^n
 \frac{\partial^2P}{\partial X_{\tau i b}\,
                 \partial X_{\bar\tau j b}}=0
 \quad\hbox{for every $\tau,i,j$.}
\end{equation}
Write $D_u$ for the invariant derivation on $G_K$ with tangent vector $u$.
Given $m\geq0$ and $f_b\in R[[x_1,\ldots,x_g]]$, define the harmonic value
\begin{equation}\label{eq:discriminant-harmonic-value}
 H(P;f)=\sum_\alpha P_\alpha
   \prod_{b=1}^n
    \left[\prod_{\tau,i}D_{u_{\tau i}}^{\alpha_{\tau i b}}
                  \bigl(e^{m h_L}f_b\bigr)\right]_{x=0}.
\end{equation}
The exponential is taken over $K$; its coefficients need not be integral.

\begin{theorem}\label{thm:discriminant-harmonic-clearing}
Under these hypotheses,
\[
                  \Delta^d H(P;f)\in R.
\]
No assumption that $2$ is a unit, that $E/\mathbb Q$ is Galois, or that
complex conjugation moves the primes above the residue characteristic is
needed.  The given integral biextension and the element $w$ remain part of
the hypotheses.
\end{theorem}

\begin{proof}
We give the integral and differential details, since they control possible
denominators at the prime $2$.

\emph{The integral module.}
Use the pairing
$\beta((u,v),(u',v'))=u^t v'-v^t u'$ on $K^g\oplus K^g$ and the lattice
\[
             \mathcal H_R=\{(u,\phi-Su):u,\phi\in R^g\}.
\]
This is a free $R$-module with the displayed coordinates; a tangent vector
$u$ has the integral lift $(u,-Su)$.  If $A_a$ is the differential of $[a]$,
the action on $K^g\oplus K^g$ is
\[
             a\cdot(u,v)=(A_a u,A_{\bar a}^{t}v).
\]
It preserves $\mathcal H_R$.  Indeed $A_a$ is integral, and the mixed
coefficient of the adjoint identification of the integral biextension gives
\begin{equation}\label{eq:discriminant-graph-integral}
                    S A_a-A_{\bar a}^{t}S\in M_g(R).
\end{equation}
The new second integral coordinate is
$A_{\bar a}^{t}\phi+(S A_a-A_{\bar a}^{t}S)u$.
For precision, the logarithm occurring in that adjoint identification has
zero constant and linear terms.  Its mixed degree-two coefficient equals
the corresponding coefficient of its exponential.  Thus
\eqref{eq:discriminant-graph-integral} is an integral coefficient of the
given unit cocycle, with no division by $2$.  Also $S_{ij}$ is the negative
coefficient of $x_jz_i$ in $h$, with no extra factor of $2$.

\emph{Integral eigenlifts after multiplication by $\Delta$.}
Choose an integral basis $b_1,\ldots,b_r$ of $O$ and its trace-dual basis
$b_1^*,\ldots,b_r^*$ in $E$.  The trace matrix has determinant
$\operatorname{disc}E$, so its adjugate shows that $\Delta b_i^*\in O$.
Every embedding of $O$ into $K$ takes values in $R$, because $R$ is
integrally closed.  The projector on the $\tau$-eigenspace is
\[
 e_\tau=\sum_i\tau(b_i^*)[b_i].
\]
The trace-duality identity gives $e_\tau$ this action on every split
$E$-module.  Consequently $\Delta e_\tau$ preserves $\mathcal H_R$.
Apply it to $(u_{\tau i},-S u_{\tau i})$.  We obtain an integral eigenlift
\[
 \xi_{\tau i}=(u'_{\tau i},v_{\tau i})\in\mathcal H_R,
                  \qquad u'_{\tau i}=\Delta u_{\tau i}.
\]
The pairing is adjoint for conjugation:
$\beta(a\xi,\eta)=\beta(\xi,\bar a\eta)$.
Applying this to the horizontal eigenvector $(u'_{\nu j},0)$ and to
$\xi_{\tau i}$ shows
\begin{equation}\label{eq:discriminant-contraction-support}
          (u'_{\nu j})^t v_{\tau i}=0
                         \quad\hbox{unless $\nu=\bar\tau$.}
\end{equation}
Distinct embeddings can be separated by an element of $O$; the resulting
nonzero scalar is cancelled in the field $K$, without asserting that it
is a unit of $R$.

\emph{Integral differential words.}
Let $\nabla_{(u,v)}$ denote the first-order operator on the formal
biextension with central weight $m$.  Put $U=e^{m h_L}$.
For $\xi=(u,\phi-Su)\in\mathcal H_R$, the conjugate operator
$U^{-1}\nabla_\xi U$ preserves $R[[x,y]]$.
Here is the coefficient check.  Its horizontal vector-field coefficients
are those of $D_u$, hence integral.  Its vertical coefficients are
\[
             D_u\sigma_i-(Su)_i+\phi_i,
\]
and its multiplication term is
\[
        m\left(D_u h_L+\sum_i\sigma_i u_i-\sum_i y_i u_i\right).
\]
The biextension identities give
$D_u\sigma_i-(Su)_i\in R[[x]]$ and
$D_u h_L+\sum_i\sigma_i u_i\in R[[x]]$.
To obtain them, differentiate the two addition identities and the symmetry
and adjoint identities of the rigidified integral biextension; use
$[w]+[\bar w]=[1]$ in the second expression.  Each logarithmic derivative
of an integral unit is its derivative multiplied by its integral inverse.
This proves these integrality assertions using ordinary derivatives and
integer multiplication only.  Repeating such operators therefore
introduces neither factorial denominators nor a denominator $2$.

\emph{Harmonicity removes all contractions.}
Choose one member of each pair $\{\tau,\bar\tau\}$ and order the operators
from that side before those from the other side.  Write
$\nabla_{\xi_\ell}=A_\ell+B_\ell$, where
$A_\ell=\nabla_{(u'_\ell,0)}$ and
$B_\ell=\nabla_{(0,v_\ell)}$.
The $A$'s commute, the $B$'s commute, and
$[B_\ell,A_{\ell'}]=m\beta((0,v_\ell),(u'_{\ell'},0))$ is a scalar.
Within each chosen side this scalar is zero by
\eqref{eq:discriminant-contraction-support}.
Moving the $B$'s to the right of each ordered word gives the usual finite
contraction expansion.  In the polynomial notation its correction is
\[
         \sum_{r\geq0}\frac{\mathcal Q^rP}{r!},\qquad
 \mathcal Q=\sum_{\tau,i,j}c_{\tau i j}
       \sum_b\frac{\partial^2}
        {\partial X_{\tau i b}\,\partial X_{\bar\tau j b}},
\]
where the scalars $c_{\tau i j}$ are the displayed commutators with the
chosen order.  One can verify this identity by expanding the generating
exponentials for the two sides: commuting them contributes the exponential
of the scalar cross commutator.  Since the commutator vanishes on each
side, there is no additional half-diagonal term.
Equation~\eqref{eq:discriminant-harmonic} gives $\mathcal QP=0$.
Thus every term with $r>0$ is exactly zero in $K$ before integrality is
considered.  There is no division of a nonzero integral coefficient by $r!$.

The remaining horizontal value is precisely
\eqref{eq:discriminant-harmonic-value} with directions $u'$.  This follows
by restricting the horizontal operators to the graph of $\sigma$; the
axis rigidification makes their constant value equal to the corresponding
invariant derivatives on $G$.
On the other hand, the ordered expression is integral: replace each
operator by $U^{-1}\nabla_\xi U$, apply the resulting words to the integral
$f_b$, and take constant coefficients.  Since $U(0)=1$, this replacement
does not alter those constants.  The coefficients of $P$ and all the
words are integral.  Finally, every monomial of $P$ has total degree $d$,
so replacing all $u$ by $\Delta u$ multiplies the value by $\Delta^d$.
\end{proof}

\paragraph{Use in the theta denominator.}
For a valuation normalized by $v/e$, $e>0$, the theorem gives the chart
bound
\[
                        v(H(P;f))/e\geq-d\,v(\Delta)/e.
\]
The original row homogeneity makes $d$ the sum of the row degrees.  Thus
the factor $\Delta^d$ already present in the proposed clearing integer
\[
                        Z=|Q|\,\Delta^d M^d
\]
absorbs this entire local loss.  Here $Q$ is the finite set of level
characters and $M$ is the level.  This replaces the sharp ramified
harmonic estimate in that application, including its pairing of ramified
primes.  It does not prove the normalized primary-torsion trace identities,
the determinant-ideal scalar comparison, theta spanning, or the independent
archimedean norm estimate.  Those inputs are still required to place the
actual theta family in the global section lattice.

The trace-discriminant calculation is standard
\cite{stacks-harmonic-discriminant}.  The formal identities used here,
including the integral graph, integral first-order operators, and the
contraction formula, are matched to the fixed original source
\cite{alpoge-harmonic-discriminant}.
